% !PRINTED 691
\weberpage{691}

%<<CONT
eine Funktion in $\Omega$, die wir einen \gesperrt{Differentialquotienten}
erster Gattung nennen.

Ist $\mathfrak{P}$ ein Punkt, in dem $(z - z_0)$ unendlich klein von der
Ordnung $e$ wird, so kommt dieser Punkt $(e - 1)$ mal in $\mathfrak{Z}$ vor,
und folglich ist $[w(z - z_0)]_0 = 0$. Wird $z$ in einem anderen
Punkte $\mathfrak{P}$ unendlich in der $e$ten Ordnung, so kommt dieser
Punkt $e$ mal in $\mathfrak{U}$ und $(e - 1)$ mal in $\mathfrak{Z}$ vor, also einmal im
Zähler von $wz$, und folglich ist auch in diesem Punkte $(wz)_0 = 0$,
und diese beiden Forderungen, nämlich:

a) In jedem Punkte, in dem $z$ einen endlichen Wert $z_0$ hat,
ist $[w(z - z_0)]_0 = 0$;

b) In einem Punkte, in dem $z$ unendlich wird, ist $(wz)_0 = 0$,
sind auch ausreichend, einen Differentialquotienten erster Gattung
$w$ zu definieren.

Um die erste Bedingung zu erfüllen, nehmen wir eine
Minimalbasis nach $z$:
\weq{2}{\omega_1,\ \omega_2,\ \dots,\ \omega_n,}
setzen
\weq{3}{a_{r,\,s} = S(\omega_r\omega_s),}
so daß
\weq{4}{\Delta = \Sigma \pm a_{1,1}a_{2,2}\dots a_{n,n}}
die Körperdiskriminante (nach $z$) ist, und bestimmen eine Basis
\weq{5}{\varepsilon_1,\ \varepsilon_2,\ \dots,\ \varepsilon_n}
aus den Gleichungen:
\weq{6}{\omega_r = \Sigma a_{r,\,i}\varepsilon_i.}

Ist $(r, s)$ ein Zeichen, das $= 0$ ist, wenn $r$ und $s$ verschieden
sind, und $= 1$, wenn $r = s$ ist, so kann man die rationalen
Funktionen $a'_{r,\,s}$ von $z$ so bestimmen, daß
\weq{7}{\sideset{^i}{}\sum a_{i,\,r}a'_{i,\,s} = (r,\ s)}
wird, und erhält aus (6):
\weq{8}{\varepsilon_r = \sideset{^i}{}\sum a'_{i,\,r}\omega_i,}
und die Größen $a'_{r,\,s}$ haben keinen anderen Nenner als $\Delta$.
Daraus ergibt sich wegen (3) und (7):
\weq{9}{\begin{aligned}
S(\varepsilon_r\omega_s) &= (r,\ s), \\
a'_{r,\,s} &= S(\varepsilon_r\varepsilon_s).
\end{aligned}}
